\section{Longitudinal and Disease-Progression Modelling}

Longitudinal data provide information about how Alzheimer’s disease (AD) develops over time by following individuals across multiple clinical visits. 
While cross-sectional studies can identify differences between clinical groups at a particular point in time, longitudinal measurements allow changes within individuals to be studied. 
Representation learning methods can use these repeated measurements to learn disease-related changes and trajectories of progression.

Yang et al.\ proposed a sequential graph autoencoder for modelling longitudinal neuroimaging data in preclinical AD~\cite{Yang2021}. The model combines relationships between brain regions with information from repeated clinical visits and separates the latent representation into time-varying and time-invariant components. 
The time-varying representation captures changes across visits, while the time-invariant representation captures characteristics that remain relatively stable over time. 
This demonstrates how autoencoder-based representations can be structured to distinguish longitudinal disease-related changes from more stable subject-specific information.

Longitudinal progression has also been modelled using recurrent variational autoencoders. 
Martí-Juan et al.\ proposed a multi-channel recurrent variational autoencoder (MC-RVAE) for modelling multimodal AD progression~\cite{MartiJuan2023}. 
The model uses a shared latent representation to capture relationships between different data modalities while recurrent components model their changes over time. 
This allows longitudinal trajectories to be reconstructed and missing measurements to be estimated, demonstrating how generative latent-variable models can represent the temporal progression of multiple AD-related measurements.

Disease progression can also be represented as a continuous trajectory within a learned latent space. 
Hong et al.\ used a variational autoencoder to learn representations of tau PET images and derive a pseudo-temporal trajectory describing the progression of tau pathology~\cite{Hong2022}. 
The resulting trajectory reflected changes in the spatial distribution of tau and was associated with clinical and biological measures of AD.\@ 
Together, these studies show that learned representations can be used not only to distinguish clinical groups, but also to describe disease-related changes over time and represent progression as a continuous process.
